% sudo apt update && sudo apt install -y texlive-full
% 使用 XeLaTeX 编译: xelatex thesis.min.tex
% CCF 中国计算机学会 A 类会议论文格式
% Lark: 基于强化学习的多信号融合自适应网络拥塞控制算法

\documentclass[10pt,twocolumn]{article}

% 中文支持
\usepackage[UTF8]{ctex}
\usepackage{fontspec}

% 数学公式
\usepackage{amsmath}
\usepackage{amssymb}

% 图表
\usepackage{graphicx}
\usepackage{float}
\usepackage{booktabs}
\usepackage{multirow}
\usepackage{array}

% 算法
\usepackage{algorithm}
\usepackage{algorithmic}

% 引用和链接
\usepackage{cite}
\usepackage{hyperref}
\usepackage{url}

% 颜色
\usepackage{xcolor}

% 页面设置
\usepackage{geometry}
\geometry{a4paper,left=1.9cm,right=1.9cm,top=2.2cm,bottom=2.2cm,columnsep=0.55cm}

% 行距
\usepackage{setspace}
\setstretch{1.08}

% 标题间距
\usepackage{titlesec}
\titlespacing*{\section}{0pt}{10pt}{5pt}
\titlespacing*{\subsection}{0pt}{7pt}{3pt}
\titlespacing*{\subsubsection}{0pt}{5pt}{2pt}
\titleformat{\section}{\normalfont\large\bfseries}{\thesection}{0.5em}{}
\titleformat{\subsection}{\normalfont\normalsize\bfseries}{\thesubsection}{0.5em}{}
\titleformat{\subsubsection}{\normalfont\small\bfseries}{\thesubsubsection}{0.5em}{}

% 紧凑列表
\usepackage{enumitem}
\setlist{nosep,leftmargin=1.2em}

% caption
\usepackage{caption}
\captionsetup{font=small,labelfont=bf,skip=4pt}

% 浮动体间距
\setlength{\textfloatsep}{8pt plus 2pt minus 2pt}
\setlength{\floatsep}{6pt plus 2pt minus 2pt}
\setlength{\intextsep}{6pt plus 2pt minus 2pt}
\setlength{\abovecaptionskip}{4pt}
\setlength{\belowcaptionskip}{2pt}

% 段落间距
\setlength{\parskip}{1.5pt plus 0.5pt minus 0.5pt}

% 算法环境字体
\usepackage{etoolbox}
\AtBeginEnvironment{algorithmic}{\small}

\begin{document}

%========================================
% 标题
%========================================

\title{\Large\textbf{Lark: 基于强化学习的多信号融合自适应网络拥塞控制算法}}

\author{
\normalsize 杭天铖 \\
\small 南京邮电大学 \\
\small 1224045520@njupt.edu.cn
}

\date{}

\maketitle

%========================================
% 摘要
%========================================

\begin{abstract}
\small
现代数据中心网络的高带宽、低延迟特性对TCP拥塞控制机制提出了严峻挑战。传统基于丢包的算法存在拥塞信号滞后、缓冲区膨胀等问题，基于延迟的算法面临RTT测量不稳定和公平性差等困难。本文提出Lark算法——一种基于强化学习的多信号融合自适应网络拥塞控制方法。其核心创新包括：（1）设计包含15个参数的完整观测空间，首次将ECN状态、拥塞控制状态机状态等协议栈信息纳入强化学习状态表示；（2）提出多信号融合拥塞检测方法，采用两级优先级机制综合分析丢包、超时和ECN信号，并通过连续递减保护机制防止窗口过度收缩；（3）设计差异化拥塞响应机制，针对不同拥塞类型采用不同的窗口保留因子（丢包0.70、ECN 0.85、超时0.50）；（4）实现强化学习驱动的自适应参数调优，基于三级RTT梯度和奖励指数移动平均动态调整乘性增加因子$\alpha \in [1.05, 1.50]$；（5）采用基于带宽延迟积窗口最大值滤波估计与当前窗口的融合控制策略。在ns-3平台上9种典型数据中心网络场景的实验表明，Lark在保持与传统最优算法（NewReno/CUBIC）相当吞吐量（差距$<$1\%）的同时，将延迟降低了70\%至81\%。

\vspace{0.3em}
\noindent\textbf{关键词：}拥塞控制；强化学习；多信号融合；显式拥塞通知；自适应参数调优；数据中心网络
\end{abstract}

%========================================
% 1. 引言
%========================================

\section{引言}

随着云计算、大数据和人工智能技术的蓬勃发展，全球数据中心网络的规模和复杂度呈现指数级增长态势。根据IDC统计，全球数据中心网络流量以年均25\%以上的速度持续增长\cite{idc2023}。TCP作为承载超过90\%数据中心流量的核心传输协议，其拥塞控制机制的性能直接影响网络资源利用效率、应用响应延迟和用户服务质量。

现代数据中心网络普遍采用Leaf-Spine或Fat-Tree拓扑结构，网络带宽经历了从千兆到10Gbps、25Gbps乃至100Gbps级别的跨越式提升，而网络延迟则从毫秒级降低到微秒级\cite{datacentertraffic}。机架内通信延迟通常在微秒级别，跨Pod通信延迟也仅在数十微秒量级。数据中心流量呈现显著的"Mice-Elephant"分布特征：超过80\%的流量是小于10KB的短流，但超过80\%的带宽被不到1\%的大流消耗。Incast（多对一）通信模式在分布式系统中十分常见，会导致瞬时的严重拥塞。这种高带宽、低延迟、动态突发的网络环境对传统拥塞控制算法提出了前所未有的挑战。

传统基于丢包的拥塞控制算法（如NewReno\cite{newreno}、CUBIC\cite{cubic}）将丢包作为拥塞信号，但丢包是拥塞的滞后指示器，通常意味着网络缓冲区已经溢出，排队延迟已累积到较高水平，产生"Bufferbloat"问题\cite{bufferbloat}。基于延迟的算法（如Vegas\cite{vegas}、BBR\cite{bbr}）尝试通过RTT变化提前感知拥塞，但RTT测量易受操作系统调度、路径抖动等因素干扰，且与基于丢包算法共存时存在严重的公平性问题\cite{bbrfairness}。

ECN机制\cite{rfc3168}允许网络设备在队列超过阈值时标记数据包而非丢弃，实现更早、更精确的拥塞通知。DCTCP\cite{dctcp}利用ECN标记比例进行细粒度调整，取得了良好效果。然而，现有算法对ECN信号的利用仍不够充分，大多采用静态参数配置，缺乏自适应能力。

近年来，强化学习在网络拥塞控制中的应用受到关注\cite{rlcc}。Aurora\cite{aurora}采用PPO算法训练深度神经网络策略，Orca\cite{orca}使用强化学习动态调整AIMD参数。然而，现有方案存在状态空间设计不完善（通常仅4-6维）、未利用ECN和拥塞控制状态机信息、奖励函数难以平衡多目标、训练效率低等问题。

现代数据中心普遍采用Leaf-Spine拓扑，具有多路径冗余、低网络直径（4-6跳）和均匀延迟等特性。数据中心流量呈显著"Mice-Elephant"分布：超过80\%的流量小于10KB，但超过80\%的带宽被不到1\%的大流消耗。Incast通信模式在分布式系统中普遍存在，会导致瞬时严重拥塞。不同应用对网络性能需求各异：在线交易处理（OLTP）要求P99延迟控制在亚毫秒级，分布式存储需要高吞吐量，机器学习训练对带宽利用率和公平性有较高要求。同时，现代数据中心交换机普遍采用浅缓冲区设计（每端口数百KB到数MB），这使得拥塞发生时更容易丢包，对算法响应速度提出了更高要求。

针对上述挑战，本文提出Lark算法——一种基于强化学习的多信号融合自适应网络拥塞控制方法。"Lark"（云雀）象征高吞吐量与低延迟的设计目标。本文主要贡献如下：

\textbf{（1）15维观测空间}：首次将ECN状态、拥塞控制状态机状态、拥塞事件类型等协议栈信息纳入强化学习状态表示，提供更完整的网络状态感知。

\textbf{（2）多信号融合拥塞检测}：采用两级优先级机制综合分析丢包、超时和ECN信号，通过连续递减保护机制防止窗口过度收缩。

\textbf{（3）差异化拥塞响应}：针对不同拥塞类型采用差异化窗口保留因子（丢包0.70、ECN 0.85、超时0.50），精准匹配拥塞严重程度。

\textbf{（4）自适应参数调优与融合控制}：基于三级RTT梯度（1.3/2.0/3.0）、奖励指数移动平均和连续增长趋势动态调整乘性增加因子$\alpha \in [1.05, 1.50]$，将基于BDP窗口最大值滤波估计的速率控制与基于当前窗口的保守控制相结合。

%========================================
% 2. 相关工作
%========================================

\section{相关工作}

\subsection{传统拥塞控制算法}

\subsubsection{基于丢包的算法}

TCP Tahoe和Reno是最早的拥塞控制算法\cite{jacobson1988}，引入了慢启动、拥塞避免和快速重传/恢复机制。TCP NewReno\cite{newreno}改进了快速恢复机制，能够处理突发丢包情况，是RFC 6582定义的标准算法。

TCP CUBIC\cite{cubic}是Linux自2.6.19版本以来的默认算法，采用三次函数模型调整窗口：
\begin{equation}
W(t) = C \cdot (t - K)^3 + W_{max}
\end{equation}
其中$K = \sqrt[3]{W_{max} \cdot \beta / C}$。CUBIC的窗口增长与RTT无关，在高BDP网络中具有更快的收敛速度。

基于丢包算法的共同缺陷包括：（1）拥塞信号滞后，缓冲区溢出前已积累大量排队延迟；（2）缓冲区膨胀，持续增加发送速率直至丢包；（3）无法区分拥塞程度，对所有拥塞采用统一的窗口减半策略。

\subsubsection{基于延迟的算法}

TCP Vegas\cite{vegas}通过比较期望吞吐量$cwnd/BaseRTT$与实际吞吐量$cwnd/RTT$的差值$\Delta$调整窗口，当$\Delta < \alpha$时增窗，$\Delta > \beta$时降窗。

BBR\cite{bbr}通过估计瓶颈带宽$BtlBw$和最小RTT $RTT_{prop}$，以BDP为目标调整发送速率：$pacing\_rate = pacing\_gain \times BtlBw$。BBR采用周期性增益调整实现带宽探测。

Copa\cite{copa}将目标排队延迟设为$RTT_{standing} - RTT_{min}$的函数，并引入竞争模式检测机制。

这类算法面临RTT测量不稳定（受调度延迟、路径抖动影响）、与丢包算法共存公平性差（研究表明BBR与CUBIC共存时吞吐量可下降30\%以上\cite{bbrfairness}）、参数敏感等问题。

\subsubsection{混合算法}

Compound TCP\cite{compound}叠加基于延迟的窗口分量，Illinois\cite{illinois}根据RTT调整AIMD参数，Westwood\cite{westwood}通过ACK到达率估计带宽。混合算法在一定程度上缓解了单一信号的局限性，但增加了复杂性。

\subsection{ECN机制与相关算法}

ECN\cite{rfc3168}通过IP报头中的2比特字段传递拥塞信息，定义四种码点（Non-ECT、ECT(0)、ECT(1)、CE）。发送端设置ECT标记表明支持ECN；网络设备在队列超阈值时将ECT改为CE；接收端回传ECE标志；发送端响应后设置CWR标志。

DCTCP\cite{dctcp}利用ECN标记比例$\alpha$进行细粒度调整：
\begin{equation}
\alpha = (1-g) \cdot \alpha + g \cdot F, \quad cwnd = cwnd \times (1 - \alpha/2)
\end{equation}
其中$g=1/16$是平滑因子，$F$是最近一个RTT内被标记的数据包比例。DCTCP能根据拥塞程度进行比例响应，但参数静态预设。HPCC\cite{hpcc}利用网络设备提供的INT精确拥塞信息进行控制，但需要特殊硬件支持。

ECN的优势在于：（1）在丢包前提供明确拥塞信号；（2）避免重传开销；（3）信号来自网络设备显式标记，比RTT推测更可靠。然而，现有算法对ECN的利用仍不充分。

\subsection{基于强化学习的拥塞控制}

Aurora\cite{aurora}采用PPO算法训练神经网络策略，使用10维状态空间（延迟梯度、延迟比率、发送比率等），奖励函数为$r = throughput - \delta \cdot delay - \lambda \cdot loss$。Orca\cite{orca}在AIMD框架下用强化学习动态调整增加因子$\alpha$和减少因子$\beta$。PCC\cite{pcc}定义实用函数$U(x) = T \cdot x^\alpha - \delta \cdot L \cdot x - \beta \cdot D \cdot x$，通过梯度上升最大化。Remy\cite{remy}采用离线优化生成规则。

现有方法的共同问题是：（1）状态空间不完善，仅4-6维，未利用ECN和CA状态信息；（2）奖励函数设计困难，多目标权衡复杂；（3）训练效率低；（4）泛化能力有限。Lark通过15维观测空间、多信号融合和差异化响应系统性地解决这些问题。

%========================================
% 3. Lark算法设计
%========================================

\section{Lark算法设计}

\subsection{系统架构}

Lark的系统架构由四个核心模块组成。\textbf{网络仿真模块}基于ns-3\cite{ns3}构建，实现完整TCP/IP协议栈和多种网络拓扑（Leaf-Spine、哑铃型等），采用MTP多线程并行化架构，通过保守同步协议实现3-5倍仿真加速。\textbf{拥塞控制环境模块}基于OpenGym\cite{opengym}框架实现，在TCP协议栈的五个关键回调点（GetSsThresh、IncreaseWindow、PktsAcked、CongestionStateSet、CwndEvent）采集15维状态参数，并将智能体决策应用于TCP连接。\textbf{智能决策模块}采用Python实现，执行多信号融合拥塞检测、差异化拥塞响应、自适应参数调优和融合控制策略计算，维护每个TCP连接的独立状态。\textbf{进程间通信模块}采用ZeroMQ消息队列和Protocol Buffers序列化，实现仿真进程与决策进程间的高效通信。

Lark将拥塞控制建模为马尔可夫决策过程$\langle S, A, P, R, \gamma \rangle$。状态空间$S$为15维观测向量，动作空间$A$直接输出新的ssthresh和cwnd值（连续动作空间），采用基于规则的确定性策略而非神经网络，保证实时性（$\sim$0.1ms决策延迟）和可解释性。

状态空间设计满足以下理论要求：（1）\textbf{马尔可夫性}——当前状态包含做出最优决策所需的全部信息，即$P(s_{t+1}, r_{t+1} | s_t, a_t, s_{t-1}, ...) = P(s_{t+1}, r_{t+1} | s_t, a_t)$，通过包含ECN状态、CA状态等协议栈内部状态近似满足；（2）\textbf{信息充分性}——15维空间涵盖TCP协议栈所有关键信息，避免因信息不足导致的次优决策；（3）\textbf{可观测性}——所有参数均可从ns-3标准回调接口直接获取；（4）\textbf{维度可控性}——15维处理时间复杂度$O(1)$，基于规则的决策避免了维度灾难问题。

\subsection{15维观测空间设计}

Lark设计了包含15个参数的完整观测空间，如表\ref{tab:observation}所示。参数涵盖六大类别：连接标识（socketUuid、nodeId）用于多流识别和状态隔离；时间信息（envType、simTime）用于时序分析和投递速率计算；窗口状态（ssthresh、cwnd）是拥塞控制的核心变量；传输性能（segSize、segmentsAcked、bytesInFlight）用于吞吐量估算和管道利用率计算；延迟测量（lastRtt、minRtt）用于BDP估算和拥塞程度判断；回调类型与拥塞控制状态（callbackType、caState、caEvent、ecnState）提供协议栈完整状态。

\begin{table}[htbp]
\centering
\caption{Lark 15维观测空间参数定义}
\label{tab:observation}
\small
\begin{tabular}{clp{3.2cm}}
\toprule
\textbf{索引} & \textbf{参数} & \textbf{说明} \\
\midrule
0 & socketUuid & 套接字唯一标识 \\
1 & envType & 环境类型(0/1) \\
2 & simTime & 仿真时间($\mu$s) \\
3 & nodeId & 节点标识 \\
4 & ssthresh & 慢启动阈值(B) \\
5 & cwnd & 拥塞窗口(B) \\
6 & segSize & 段大小(B) \\
7 & segmentsAcked & 确认段数 \\
8 & bytesInFlight & 在途字节数 \\
9 & lastRtt & 最近RTT($\mu$s) \\
10 & minRtt & 最小RTT($\mu$s) \\
11 & callbackType & 回调类型(0/1) \\
12 & caState & CA状态(0-4) \\
13 & caEvent & CA事件(0-5) \\
14 & ecnState & ECN状态(0-5) \\
\bottomrule
\end{tabular}
\end{table}

Lark的核心创新在于引入三个协议栈内部状态参数。\textbf{caState}反映TCP拥塞控制状态机的当前状态：CA\_OPEN（0，正常传输）、CA\_DISORDER（1，收到重复确认）、CA\_CWR（2，响应ECN缩减窗口）、CA\_RECOVERY（3，快速恢复中）、CA\_LOSS（4，超时丢包）。\textbf{caEvent}表示最近的拥塞事件：包括CA\_EVENT\_TX\_START（0）、CWND\_RESTART（1）、COMPLETE\_CWR（2）、LOSS（3）、ECN\_NO\_CE（4）、ECN\_IS\_CE（5）六种。\textbf{ecnState}反映ECN协商和标记状态：ECN\_DISABLED（0）、ECN\_IDLE（1）、ECN\_CE\_RCVD（2）、ECN\_SENDING\_ECE（3）、ECN\_ECE\_RCVD（4）、ECN\_CWR\_SENT（5）。

这三个参数的引入使智能体能够感知TCP协议栈的完整状态信息，近似满足马尔可夫性要求。所有15个参数均可从ns-3 TCP模块的标准回调接口直接获取，满足可观测性要求。15维的处理时间复杂度为$O(1)$，不影响决策延迟。

\subsection{多信号融合拥塞检测}

Lark将拥塞信号分为三类，采用两级优先级层次结构进行综合判断，并引入连续递减保护机制防止窗口过度收缩。

\textbf{第一优先级——丢包与超时检测}：当$callbackType = 0$时（GetSsThresh回调），表明协议栈检测到拥塞事件。此时进一步区分超时与普通丢包：若$caState = \text{CA\_LOSS}$，判定为超时丢包（通常意味着连续多个数据包丢失或网络路径中断）；否则判定为普通丢包（快速重传触发）。丢包和超时是最可靠的拥塞信号（似然比$LR \approx 100$），一旦检测到即触发响应并重置连续递减计数器。

\textbf{第二优先级——ECN状态直接检测}：当$ecnState \in \{\text{ECN\_CE\_RCVD}, \text{ECN\_ECE\_RCVD}\}$时，判定存在ECN拥塞标记。与基于事件率统计的方法不同，Lark直接检查协议栈的ECN状态字段，避免了维护时间戳队列的额外开销和事件率阈值的参数敏感性问题。

\textbf{连续递减保护机制}：为防止在突发拥塞期间窗口被过度收缩，Lark维护一个连续递减计数器$d$。每次触发拥塞响应时$d$自增；每次正常窗口增长时$d$重置为零。当$d$超过最大连续递减次数$D_{max} = 3$时，拥塞响应被抑制——保持当前窗口不变，避免因连续乘性递减导致窗口收缩至过小值。理论分析表明，连续3次丢包响应后窗口仅剩$0.70^3 \approx 34.3\%$，已接近传统算法单次窗口减半的效果，进一步缩减的边际收益递减。

完整的多信号融合拥塞检测算法如算法\ref{alg:detect}所示。

\begin{algorithm}[htbp]
\caption{Multi-Signal Fusion Congestion Detection}
\label{alg:detect}
\small
\begin{algorithmic}[1]
\REQUIRE Observation $obs$, consecutive decrease counter $d$
\ENSURE $(is\_congested, \; type)$
\STATE \textbf{// Level-1: Packet loss and timeout}
\IF{$obs.callbackType = 0$}
    \STATE $d \leftarrow 0$
    \IF{$obs.caState = \text{CA\_LOSS}$}
        \RETURN $(True, \; \text{TIMEOUT})$
    \ELSE
        \RETURN $(True, \; \text{LOSS})$
    \ENDIF
\ENDIF
\STATE \textbf{// Level-2: ECN state inspection}
\IF{$obs.ecnState \in \{\text{CE\_RCVD}, \text{ECE\_RCVD}\}$}
    \RETURN $(True, \; \text{ECN})$
\ENDIF
\RETURN $(False, \; \text{NONE})$
\end{algorithmic}
\end{algorithm}

\subsection{差异化拥塞响应机制}

传统算法（如NewReno、CUBIC）对所有拥塞信号采用统一的窗口减半策略（保留因子$\rho=0.5$），这种"一刀切"的方法无法区分轻度和严重拥塞，在轻度拥塞时响应过度（不必要地降低吞吐量），在严重拥塞时响应不足。Lark根据拥塞类型设计差异化的窗口保留因子：

\begin{equation}
new\_cwnd = \rho \cdot cwnd, \quad \rho = \begin{cases}
0.70 & \text{丢包(LOSS)} \\
0.85 & \text{ECN标记} \\
0.50 & \text{超时(TIMEOUT)}
\end{cases}
\end{equation}

设计依据如下。（1）\textbf{拥塞信号的时序特性}：ECN标记是最早的信号（队列超低阈值时触发），丢包较晚（队列溢出），超时最晚（连续丢包或路径问题）。早期信号应采用温和响应（高保留因子），晚期信号采用激进响应（低保留因子）。（2）\textbf{信号可靠性}：丢包几乎不会误报（$LR \approx 100$），ECN的可靠性取决于设备配置，因此对丢包响应更坚决。（3）\textbf{吞吐量-延迟权衡}：丢包保留因子0.70高于传统的0.5，配合ECN的提前预警机制实现更平滑的响应；ECN保留因子0.85仅缩减15\%窗口，足以缓解轻度拥塞同时维持高带宽利用率；超时保留因子0.50与传统窗口减半一致，反映超时信号所指示的严重拥塞或路径故障。

慢启动阈值更新为$new\_ssthresh = \max(new\_cwnd, BDP)$，以BDP为下界加速恢复。窗口下限约束为$4 \times MSS$。差异化响应算法如算法\ref{alg:response}所示。

\begin{algorithm}[htbp]
\caption{Differentiated Congestion Response}
\label{alg:response}
\small
\begin{algorithmic}[1]
\REQUIRE Current $cwnd$, congestion type $type$, consecutive decrease counter $d$, BDP estimate $bdp$
\ENSURE New window $new\_cwnd$, new threshold $new\_ssthresh$
\STATE $d \leftarrow d + 1$
\IF{$d > D_{max}$}
    \STATE $new\_cwnd \leftarrow \max(cwnd, \; 4 \times MSS)$ \COMMENT{Freeze window}
    \STATE $new\_ssthresh \leftarrow new\_cwnd$
    \RETURN $(new\_cwnd, \; new\_ssthresh)$
\ENDIF
\IF{$type = \text{LOSS}$}
    \STATE $\rho \leftarrow 0.70$
\ELSIF{$type = \text{ECN}$}
    \STATE $\rho \leftarrow 0.85$
\ELSIF{$type = \text{TIMEOUT}$}
    \STATE $\rho \leftarrow 0.50$
\ENDIF
\STATE $new\_cwnd \leftarrow \max(\rho \times cwnd, \; 4 \times MSS)$
\STATE $new\_ssthresh \leftarrow \max(new\_cwnd, \; bdp)$
\RETURN $(new\_cwnd, \; new\_ssthresh)$
\end{algorithmic}
\end{algorithm}

\subsection{强化学习驱动的自适应参数调优}

Lark根据多种网络状态因素动态调整乘性增加因子$\alpha \in [1.05, 1.50]$（初始值1.20），控制拥塞避免阶段窗口增长速度。调优基于三个维度的反馈信号。

\textbf{基于三级RTT梯度的调整}：RTT膨胀比$r = lastRtt/minRtt$反映排队延迟程度。Lark采用三级梯度提供细粒度控制：当$r < 1.3$时，网络几乎无排队延迟，$\alpha \mathrel{+}= 0.03$加速窗口增长并递增连续增长计数器；当$1.3 \leq r < 2.0$时，存在中等排队延迟，$\alpha \mathrel{+}= 0.01$温和增长；当$r > 3.0$时，排队延迟严重，$\alpha \mathrel{-}= 0.02$较强减速并重置连续增长计数器。

\textbf{基于奖励指数移动平均（EMA）的调整}：为捕获长期性能趋势，维护奖励信号的指数移动平均$\bar{r}_t = (1 - \eta) \cdot \bar{r}_{t-1} + \eta \cdot r_t$，其中平滑因子$\eta = 0.1$。当$\bar{r} > 2.0$时网络性能持续向好，$\alpha \mathrel{+}= 0.01$；当$\bar{r} < -5.0$时性能持续恶化，$\alpha \mathrel{-}= 0.01$。

\textbf{基于连续增长趋势的调整}：窗口连续增长超过5次时，表明网络状况持续良好，$\alpha \mathrel{+}= 0.02$并重置计数器，实现"渐进式加速"。

参数调整采用小步长（0.01$\sim$0.03），边界约束$\alpha \in [1.05, 1.50]$保证稳定性。下界1.05保证基本窗口增长率，上界1.50防止过度激进。从控制论角度，$\alpha$调整等价于根据三级延迟反馈$\tau$（RTT）动态调整控制增益$K$，满足稳定性条件$K\tau < \pi/2$\cite{tcp_congestion}。自适应参数调优算法如算法\ref{alg:adapt}所示。

\begin{algorithm}[htbp]
\caption{RL-Driven Adaptive Parameter Tuning}
\label{alg:adapt}
\small
\begin{algorithmic}[1]
\REQUIRE Current $\alpha$, observation $obs$, consecutive increase count $g$, reward EMA $\bar{r}$
\ENSURE Updated $\alpha$
\STATE $r \leftarrow obs.lastRtt \;/\; obs.minRtt$ \COMMENT{RTT inflation ratio}
\STATE \textbf{// Three-level RTT gradient}
\IF{$r < 1.3$}
    \STATE $\alpha \leftarrow \alpha + 0.03$; $g \leftarrow g + 1$
\ELSIF{$r < 2.0$}
    \STATE $\alpha \leftarrow \alpha + 0.01$
\ELSIF{$r > 3.0$}
    \STATE $\alpha \leftarrow \alpha - 0.02$; $g \leftarrow 0$
\ENDIF
\STATE \textbf{// Reward EMA trend}
\IF{$\bar{r} > 2.0$}
    \STATE $\alpha \leftarrow \alpha + 0.01$
\ELSIF{$\bar{r} < -5.0$}
    \STATE $\alpha \leftarrow \alpha - 0.01$
\ENDIF
\STATE \textbf{// Consecutive growth acceleration}
\IF{$g > 5$}
    \STATE $\alpha \leftarrow \alpha + 0.02$; $g \leftarrow 0$
\ENDIF
\STATE $\alpha \leftarrow \text{clamp}(\alpha, \; 1.05, \; 1.50)$
\end{algorithmic}
\end{algorithm}

\subsection{融合控制策略}

Lark将基于BDP的速率控制与基于当前窗口的保守控制相结合。BDP估计采用窗口最大值滤波器（Windowed Max-Filter）：投递速率$DR = segmentsAcked \times segSize / lastRtt$存入容量为$W_{bw}=20$的滑动窗口队列，取窗口内最大值$BW_{max}$作为瓶颈带宽估计；最小RTT持续跟踪全局最小值$RTT_{min}$。BDP计算为$BDP = BW_{max} \times RTT_{min}$。当BDP估计无效时回退至当前窗口值。BDP估计算法如算法\ref{alg:bdp}所示。

\begin{algorithm}[htbp]
\caption{Windowed Max-Filter BDP Estimation}
\label{alg:bdp}
\small
\begin{algorithmic}[1]
\REQUIRE Flow state $state$, observation $obs$
\ENSURE BDP estimate
\STATE $dr \leftarrow obs.segAcked \times obs.segSize \;/\; obs.lastRtt$
\STATE $state.bw\_samples$.append($dr$) \COMMENT{Sliding window, capacity $W_{bw}=20$}
\STATE $BW_{max} \leftarrow \max(state.bw\_samples)$
\STATE $RTT_{min} \leftarrow \min(state.min\_rtt, \; obs.minRtt)$
\STATE $bdp \leftarrow BW_{max} \times RTT_{min}$
\IF{$bdp \leq 0$}
    \STATE $bdp \leftarrow \max(obs.cwnd, \; 1)$ \COMMENT{Fallback}
\ENDIF
\RETURN $bdp$
\end{algorithmic}
\end{algorithm}

\textbf{慢启动阶段}（$cwnd < ssthresh$）：目标窗口$target = 2 \times BDP$（下限$10 \times MSS$）。增量$\Delta = segmentsAcked \times MSS$（标准指数增长）；当$cwnd < 0.3 \times BDP$时增量翻倍至$2 \times segmentsAcked \times MSS$，实现快速带宽探测。到达目标后自动退出慢启动阶段。

\textbf{拥塞避免阶段}（$cwnd \geq ssthresh$）：
\begin{equation}
target = \max(\alpha \times BDP, \; cwnd) + \gamma \times MSS
\end{equation}
其中$\gamma=2.0$提供恒定的加性增长偏置。窗口直接设为目标值，实现BDP导向的精确控制。

\textbf{管道利用率感知加速}：管道利用率$u = bytesInFlight / cwnd$反映网络管道的充填程度。当$u < 0.7$时，额外增加$1 \times MSS$；当$u < 0.4$时，再额外增加$2 \times MSS$（两个条件均满足时累计$3 \times MSS$），确保管道未充分利用时快速适应带宽变化。

窗口约束为$cwnd \in [4 \times MSS, \max(10 \times BDP, 200 \times MSS)]$。融合控制策略的完整算法如算法\ref{alg:fusion}所示。

\begin{algorithm}[htbp]
\caption{Hybrid Rate-Window Fusion Control}
\label{alg:fusion}
\small
\begin{algorithmic}[1]
\REQUIRE Observation $obs$, parameters $\alpha, \gamma$, BDP estimate $bdp$
\ENSURE New window $new\_cwnd$
\STATE $d \leftarrow 0$ \COMMENT{Reset consecutive decrease counter}
\STATE $u \leftarrow obs.bytesInFlight \;/\; obs.cwnd$
\IF{$obs.cwnd < obs.ssthresh$}
    \STATE $target \leftarrow \max(2 \times bdp, \; 10 \times MSS)$
    \STATE $\Delta \leftarrow segAcked \times MSS$
    \IF{$obs.cwnd < 0.3 \times bdp$}
        \STATE $\Delta \leftarrow 2 \times segAcked \times MSS$
    \ENDIF
    \STATE $new\_cwnd \leftarrow \min(obs.cwnd + \Delta, \; target)$
\ELSE
    \STATE $new\_cwnd \leftarrow \max(\alpha \times bdp, \; obs.cwnd) + \gamma \times MSS$
\ENDIF
\IF{$u < 0.4$}
    \STATE $new\_cwnd \leftarrow new\_cwnd + 3 \times MSS$
\ELSIF{$u < 0.7$}
    \STATE $new\_cwnd \leftarrow new\_cwnd + 1 \times MSS$
\ENDIF
\STATE $new\_cwnd \leftarrow \text{clamp}(new\_cwnd, \; 4 \times MSS, \; \max(10 \times bdp, \; 200 \times MSS))$
\RETURN $new\_cwnd$
\end{algorithmic}
\end{algorithm}

\subsection{每流独立状态管理}

Lark为每个TCP连接维护独立状态，通过socketUuid为键的哈希表实现$O(1)$查找。每流状态包括四个子模块：（1）\textbf{带宽估计模块}——容量$W_{bw}=20$的投递速率滑动窗口队列和容量$W_{rtt}=40$的RTT采样队列，支持窗口最大值带宽估计和全局最小RTT追踪；（2）\textbf{拥塞控制计数器}——连续递减计数器$d$（用于递减保护）、连续增长计数器$g$（用于加速触发）、累计丢包和ECN事件计数；（3）\textbf{性能指标聚合}——吞吐量指数移动平均、当前BDP估计、上一步窗口和时间戳；（4）\textbf{自适应参数}——$\alpha$（初始1.20）、$\gamma$（默认2.0）和慢启动标志位。

状态管理遵循延迟初始化（首次访问时创建）和增量更新（仅更新变化字段）两个原则。每流独立管理避免了多流并发场景下的流间干扰，支持差异化控制。全局共享奖励EMA（平滑因子$\eta=0.1$）捕获跨流的整体性能趋势。

%========================================
% 4. 实验评估
%========================================

\section{实验评估}

\subsection{实验配置}

\textbf{硬件与软件}：Intel i7-10700K 8核CPU，32GB DDR4内存，Ubuntu 20.04，ns-3.40\cite{ns3}，OpenGym\cite{opengym}，Python 3.8。

\textbf{网络拓扑}：经典哑铃型拓扑——3个发送端经接入链路连接路由器R0，R0通过瓶颈链路（DropTail/RED队列）连接R1，R1经接入链路连接3个接收端。使用RED队列，标记阈值为队列长度的30\%。

\textbf{仿真参数}：仿真时长20秒，MSS=1460字节，初始窗口10个MSS（RFC 6928），启用ECN和SACK，随机种子42，每配置重复10次取平均。

\textbf{对比算法}：NewReno\cite{newreno}（经典AIMD）、CUBIC\cite{cubic}（Linux默认）、BBR\cite{bbr}（Google基于延迟）。

\textbf{评测指标}：吞吐量$T = TotalBytes / SimTime$（Mbps）、平均延迟（ms）、Jain公平性指数$J = (\sum x_i)^2 / (n \cdot \sum x_i^2)$\cite{fairness}、丢包率。

\subsection{实验场景设计}

为全面评估Lark在不同网络环境下的性能，设计了9种典型场景，涵盖机架内高速网络、不同过载比、跨Pod网络和跨数据中心WAN，如表\ref{tab:scenarios}所示。

\begin{table}[htbp]
\centering
\caption{9种实验场景配置}
\label{tab:scenarios}
\small
\begin{tabular}{clccc}
\toprule
\textbf{编号} & \textbf{场景描述} & \textbf{瓶颈} & \textbf{延迟} & \textbf{过载比} \\
\midrule
1 & 机架内高速 & 10G & 2$\mu$s & 2.5:1 \\
2 & 标准过载 & 2.5G & 5$\mu$s & 4:1 \\
3 & 高带宽过载 & 10G & 5$\mu$s & 4:1 \\
4 & 中度拥塞 & 2G & 5$\mu$s & 5:1 \\
5 & 重度拥塞 & 1G & 5$\mu$s & 10:1 \\
6 & 跨Pod 10G & 10G & 50$\mu$s & 2.5:1 \\
7 & 跨Pod 20G & 20G & 50$\mu$s & 2:1 \\
8 & 跨DC WAN & 1G & 5ms & 10:1 \\
9 & 对称低带宽 & 1G & 20$\mu$s & 1:1 \\
\bottomrule
\end{tabular}
\end{table}

\subsection{主要实验结果}

表\ref{tab:throughput}和表\ref{tab:latency}分别汇总了所有场景的吞吐量和延迟对比。

\begin{table}[htbp]
\centering
\caption{吞吐量综合对比（Mbps）}
\label{tab:throughput}
\small
\begin{tabular}{lrrrr}
\toprule
\textbf{场景} & \textbf{Lark} & \textbf{NewReno} & \textbf{CUBIC} & \textbf{BBR} \\
\midrule
场景1 & 10193 & 10213 & 10213 & 2 \\
场景2 & 2552 & 2553 & 2553 & 1299 \\
场景3 & 9931 & 10213 & 10213 & 7107 \\
场景4 & 2042 & 2042 & 2042 & 1037 \\
场景5 & 1021 & 1021 & 1021 & 522 \\
场景6 & 10109 & 10213 & 10211 & 600 \\
场景7 & 19856 & 20408 & 20404 & 604 \\
场景8 & 1012 & 1016 & 1017 & 1019 \\
场景9 & 1018 & 1021 & 1021 & 517 \\
\bottomrule
\end{tabular}
\end{table}

\begin{table}[htbp]
\centering
\caption{延迟综合对比（ms）及Lark相对NewReno的延迟降低幅度}
\label{tab:latency}
\small
\begin{tabular}{lrrrrc}
\toprule
\textbf{场景} & \textbf{Lark} & \textbf{NR} & \textbf{CUBIC} & \textbf{BBR} & \textbf{降低} \\
\midrule
场景1 & 0.458 & 2.304 & 2.294 & 0.005 & 80.1\% \\
场景2 & 0.527 & 2.325 & 2.445 & 0.018 & 77.3\% \\
场景3 & 0.436 & 2.304 & 2.298 & 0.010 & 81.1\% \\
场景4 & 0.612 & 2.405 & 2.491 & 0.036 & 74.6\% \\
场景5 & 0.837 & 2.799 & 2.755 & 0.069 & 70.1\% \\
场景6 & 0.487 & 2.309 & 2.398 & 0.061 & 78.9\% \\
场景7 & 0.512 & 1.896 & 2.207 & 0.061 & 73.0\% \\
场景8 & 5.824 & 6.739 & 7.293 & 9.081 & 13.6\% \\
场景9 & 0.743 & 2.827 & 2.777 & 0.071 & 73.7\% \\
\bottomrule
\end{tabular}
\end{table}

\textbf{吞吐量分析}：Lark在全部9个场景的吞吐量均与传统最优算法相当，差距不超过3\%，在7个场景中差距小于1\%。带宽利用率均超过99\%（场景3和7为99.3\%，其余超过100\%，含协议头部开销）。BBR在数据中心多流竞争环境下表现不佳：在数据中心内部场景吞吐量通常仅为瓶颈带宽的50\%$\sim$70\%，在跨Pod场景更是降至3\%$\sim$6\%，在机架内场景几乎完全失效（仅2Mbps）。

\textbf{延迟分析}：Lark在数据中心内部场景（场景1$\sim$7、9）将延迟从传统算法的1.9$\sim$2.8ms降至0.4$\sim$0.8ms，降幅70.1\%$\sim$81.1\%。场景3延迟降低最显著（81.1\%），场景5降幅最小但仍达70.1\%。在WAN场景（场景8）延迟降低13.6\%，因传播延迟（5ms）占主导地位，排队延迟优化空间有限。BBR虽在部分场景延迟极低（如场景1仅0.005ms），但以严重牺牲吞吐量为代价（$<$52\%带宽利用率），不具备实用价值。

\textbf{带宽利用率分析}：带宽利用率定义为$U = Throughput / Bandwidth \times 100\%$，如表\ref{tab:utilization}所示。Lark在全部9个场景中带宽利用率均超过99\%，与NewReno/CUBIC接近。BBR在数据中心内部场景仅为50\%$\sim$70\%，在跨Pod场景更降至3\%$\sim$6\%，机架内场景几乎完全失效（0.02\%）。利用率超过100\%是因计算包含了TCP/IP协议头部开销。

\begin{table}[htbp]
\centering
\caption{带宽利用率对比（\%）}
\label{tab:utilization}
\small
\begin{tabular}{lcccc}
\toprule
\textbf{场景} & \textbf{Lark} & \textbf{NR} & \textbf{CUBIC} & \textbf{BBR} \\
\midrule
场景1 (10G) & 101.9 & 102.1 & 102.1 & 0.02 \\
场景2 (2.5G) & 102.1 & 102.1 & 102.1 & 52.0 \\
场景3 (10G) & 99.3 & 102.1 & 102.1 & 71.1 \\
场景4 (2G) & 102.1 & 102.1 & 102.1 & 51.8 \\
场景5 (1G) & 102.1 & 102.1 & 102.1 & 52.2 \\
场景6 (10G) & 101.1 & 102.1 & 102.1 & 6.0 \\
场景7 (20G) & 99.3 & 102.0 & 102.0 & 3.0 \\
场景8 (1G) & 101.2 & 101.6 & 101.7 & 101.9 \\
场景9 (1G) & 101.8 & 102.1 & 102.1 & 51.7 \\
\bottomrule
\end{tabular}
\end{table}

\subsection{深入分析}

\subsubsection{收敛特性}

表\ref{tab:convergence}展示了场景2中各算法的收敛特性。Lark的收敛速度最快（0.8秒达90\%稳定吞吐量，1.2秒达99\%），稳态波动最小（$\pm$3\%），得益于：（1）融合控制策略中基于BDP窗口最大值滤波估计的目标窗口设置提供了准确的初始目标；（2）慢启动阶段的$2 \times BDP$目标快速探测带宽；（3）三级RTT梯度驱动的自适应$\alpha$在网络空闲时快速增大，加速窗口增长。NewReno最慢（2.5秒/4.0秒），因AIMD线性增长速度受限。BBR稳态波动最大（$\pm$15\%），因周期性带宽探测机制。

\begin{table}[htbp]
\centering
\caption{收敛特性对比（场景2）}
\label{tab:convergence}
\small
\begin{tabular}{lccl}
\toprule
\textbf{算法} & \textbf{90\%收敛} & \textbf{99\%收敛} & \textbf{稳态波动} \\
\midrule
Lark & 0.8s & 1.2s & $\pm$3\% \\
NewReno & 2.5s & 4.0s & $\pm$8\% \\
CUBIC & 1.5s & 2.5s & $\pm$6\% \\
BBR & 1.0s & 1.8s & $\pm$15\% \\
\bottomrule
\end{tabular}
\end{table}

\subsubsection{公平性}

在10个并发流的公平性测试中（瓶颈2.5Gbps，5$\mu$s延迟），Lark的Jain公平性指数为0.98（最大/最小流吞吐量比1.11），略低于NewReno（0.99，比值1.05），但优于CUBIC（0.97，比值1.19）和BBR（0.89，比值2.27）。Lark的每流独立状态管理和差异化响应一致性保证了良好的公平性。BBR的严重不公平源于其带宽估计在多流竞争中的不稳定性。

\subsubsection{队列行为与ECN利用}

在场景2中，各算法的瓶颈链路平均队列长度分别为：NewReno/CUBIC约120个数据包，BBR约15个，Lark仅约8个。队列长度直接决定排队延迟：$Queuing\_Delay = Queue\_Length \times Packet\_Size / Bandwidth$。对于2.5Gbps瓶颈，120个数据包对应约0.576ms排队延迟，而8个数据包仅约0.038ms。

ECN利用统计（场景2，20秒）：Lark触发ECN标记1245个（vs. NewReno 3567、CUBIC 3892），丢包仅12个（vs. NewReno 156），拥塞响应89次但每次窗口仅缩减15\%。Lark通过及时的适度ECN响应避免了队列溢出，实现了"预防胜于治疗"的拥塞控制策略。

\subsubsection{混合流量性能}

在长流（3个持续20秒）与短流（每秒5个持续0.1秒）混合场景下（瓶颈2Gbps），结果如表\ref{tab:mixed}所示。

\begin{table}[htbp]
\centering
\caption{混合流量测试结果}
\label{tab:mixed}
\small
\begin{tabular}{p{2.5cm}rrrr}
\toprule
\textbf{指标} & \textbf{Lark} & \textbf{NR} & \textbf{CUBIC} & \textbf{BBR} \\
\midrule
长流吞吐(Mbps) & 613 & 626 & 628 & 412 \\
短流完成(ms) & 12.5 & 45.3 & 42.1 & 15.8 \\
短流P99(ms) & 28.3 & 125.6 & 118.2 & 35.2 \\
链路利用率(\%) & 98.2 & 99.1 & 99.3 & 72.5 \\
\bottomrule
\end{tabular}
\end{table}

Lark的长流吞吐量略低于NewReno/CUBIC（差距$<$3\%），但短流完成时间降低72\%（12.5ms vs. 45.3ms），P99完成时间降低77\%（28.3ms vs. 125.6ms）。这对于数据中心中常见的延迟敏感短查询（如分布式数据库读取、微服务调用）具有重要实用价值。

\subsubsection{参数敏感性}

对场景2中关键参数进行敏感性分析。\textbf{丢包保留因子}：从0.50到0.90，默认值0.70在吞吐量（2552Mbps）和延迟（0.527ms）间取得最佳平衡。0.50导致吞吐量下降2\%但延迟改善13\%；0.90提升吞吐量0.8\%但延迟恶化90\%且丢包数增至35个。\textbf{ECN保留因子}：默认0.85优于0.70（吞吐量降2.5\%）和0.95（延迟增19\%）。\textbf{连续递减保护阈值}：默认$D_{max}=3$有效平衡响应灵敏度与窗口稳定性。$D_{max}=1$导致过度保守（延迟增40\%），$D_{max}=10$保护过弱（与无保护接近）。\textbf{RTT膨胀比三级梯度阈值}：默认1.3/2.0/3.0使$\alpha$均值稳定在1.20附近，三级梯度相比两级（1.5/3.0）在中等拥塞场景延迟降低约8\%。

%========================================
% 5. 讨论
%========================================

\section{讨论}

\subsection{理论分析}

\textbf{收敛性}：Lark的窗口动态可用Lyapunov函数$V = \frac{1}{2}(cwnd - cwnd^*)^2$分析，其中$cwnd^* \approx BDP$。在拥塞避免阶段，窗口向$\max(\alpha \times BDP, cwnd) + \gamma \times MSS$收敛。当$cwnd > BDP$触发ECN时，$cwnd_{t+1} = 0.85 \cdot cwnd_t < cwnd_t$，保证$\Delta V < 0$。连续递减保护机制（$D_{max}=3$）防止窗口过度收缩导致的振荡。$\alpha$的边界约束$[1.05, 1.50]$和小步长调整保证参数不发散。

\textbf{复杂度与计算开销}：单次决策时间$O(1)$，空间$O(n \cdot H)$（$n$为并发流数，$H=100$为缓冲区大小）。表\ref{tab:overhead}展示了不同硬件平台的决策延迟和内存开销。在现代x86服务器上决策延迟在0.1ms以下，远小于典型RTT；ARM平台略高但仍在1ms以下。每流状态占用5.6KB，10万并发流总内存仅$\sim$568MB。CPU利用率方面，1万次/秒决策频率仅占单核8.5\%，可支撑$\sim$100Gbps等效带宽。

\begin{table}[htbp]
\centering
\caption{计算开销分析}
\label{tab:overhead}
\small
\begin{tabular}{p{2.6cm}cc}
\toprule
\textbf{指标} & \textbf{值} & \textbf{说明} \\
\midrule
\multicolumn{3}{l}{\textit{决策延迟（均值/P99）}} \\
i7-10700K & 0.08/0.15ms & 桌面CPU \\
Xeon Gold & 0.12/0.22ms & 服务器 \\
ARM A72 & 0.35/0.65ms & 嵌入式 \\
\midrule
\multicolumn{3}{l}{\textit{内存占用}} \\
基础内存 & 12.5MB & 运行时 \\
每流状态 & 5.6KB & 缓冲+跟踪 \\
10K流 & 68.1MB & 典型规模 \\
100K流 & 568.1MB & 极端场景 \\
\midrule
\multicolumn{3}{l}{\textit{CPU利用率（单核）}} \\
1K决策/秒 & 1.2\% & $\sim$10G \\
10K决策/秒 & 8.5\% & $\sim$100G \\
\bottomrule
\end{tabular}
\end{table}

\textbf{信息论视角}：多信号融合的两级优先级等价于按似然比排序。设网络拥塞状态$C \in \{0,1\}$，各信号条件独立，则联合似然比$LR(\mathbf{S}) = \prod_i P(S_i|C=1)/P(S_i|C=0)$。丢包/超时$LR \approx 100$，ECN $LR \approx 20$。优先使用高似然比信号可最小化贝叶斯误检概率$P_e$。连续递减保护机制从信息论角度可理解为对连续同类信号的置信度递减——第$k$次连续信号的增量信息量$I_k \propto 1/k$，当$k > D_{max}$时信息增益不足以证明进一步缩减的合理性。

\textbf{最优控制视角}：差异化保留因子可从最优控制理论分析。严重拥塞（超时，连续丢包或路径故障）需最大窗口缩减（$\rho=0.50$）快速清空队列；中度拥塞（丢包，队列溢出）采用中等缩减（$\rho=0.70$）；轻度拥塞（ECN，队列开始积累）仅需适度调整（$\rho=0.85$）。Lark的三级差异化设计近似了不同拥塞程度下的最优控制输入，响应强度与拥塞严重程度正相关。

\subsection{与相关工作的详细对比}

\begin{table}[htbp]
\centering
\caption{Lark与主要相关算法的特性对比}
\label{tab:comparison}
\small
\begin{tabular}{lccc}
\toprule
\textbf{特性} & \textbf{Lark} & \textbf{Aurora} & \textbf{DCTCP} \\
\midrule
状态空间维度 & 15 & 10 & N/A \\
ECN状态利用 & 是 & 否 & 是 \\
CA状态利用 & 是 & 否 & 否 \\
决策机制 & 规则 & DNN & 公式 \\
可解释性 & 高 & 低 & 高 \\
训练需求 & 无 & 小时/天 & 无 \\
决策延迟 & 0.1ms & 1-10ms & $<$0.1ms \\
差异化响应 & 是 & 否 & 否 \\
参数自适应 & 是 & 固定 & 固定 \\
\bottomrule
\end{tabular}
\end{table}

\textbf{vs. Aurora\cite{aurora}}：Lark采用15维状态空间（vs. 10维），首次引入ECN和CA状态参数。采用基于规则决策（vs. 深度神经网络），决策延迟$\sim$0.1ms（vs. 1$\sim$10ms），无需训练过程即可部署，可解释性强便于调试优化。Aurora的优势在于神经网络的强表示能力，但在实际部署中面临训练数据收集、模型更新和推理开销等挑战。

\textbf{vs. DCTCP\cite{dctcp}}：Lark直接检测ECN协议栈状态进行离散拥塞判断（vs. 标记比例$\alpha$连续调整），实现差异化响应（丢包0.70/ECN 0.85/超时0.50 vs. 统一公式$1-\alpha/2$），综合利用丢包、ECN和CA状态三类信号（vs. 仅ECN），并实现参数自适应（vs. 静态配置）。连续递减保护机制进一步增强了在突发拥塞下的鲁棒性。

\textbf{vs. BBR\cite{bbr}}：Lark综合多种显式拥塞信号（vs. 仅RTT测量），带宽利用率$>$99\%（vs. 数据中心场景50\%$\sim$70\%），多流公平性显著更优（Jain 0.98 vs. 0.89）。BBR的带宽估计机制在多流竞争环境下存在明显缺陷，实验中在机架内和跨Pod场景严重退化。

\subsection{局限性与未来方向}

\textbf{WAN场景优化空间有限}：延迟降低仅13.6\%（vs. 数据中心内部70\%$\sim$81\%），因5ms传播延迟不可压缩。未来可通过排队延迟分离技术、多时间尺度BDP估计和自适应步长进行优化。

\textbf{参数预设}：保留因子（0.70/0.85/0.50）和阈值（连续递减保护$D_{max}=3$、RTT梯度1.3/2.0/3.0）为预设值。可通过元学习\cite{meta_learning}或贝叶斯优化在线搜索最优参数组合。

\textbf{深度RL增强}：当前基于规则的策略可解释性强但表示能力有限。可探索离线强化学习训练神经网络策略，再通过模型蒸馏提取为规则，兼顾表示能力和实时性。

\textbf{真实部署}：需将算法实现为Linux内核TCP模块并在真实数据中心测试床验证。QUIC协议的用户态实现可简化部署流程，且QUIC的多流特性可实现更精细的流间调度。

\textbf{多目标优化}：可将公平性、能耗效率纳入优化目标。帕累托最优理论为多目标权衡提供理论基础\cite{fairness}。

%========================================
% 6. 结论
%========================================

\section{结论}

本文提出了Lark——一种基于强化学习的多信号融合自适应网络拥塞控制算法。通过设计包含ECN状态、拥塞控制状态机等15维观测空间，Lark首次将协议栈内部状态信息纳入强化学习状态表示。多信号融合拥塞检测采用两级优先级机制（丢包/超时$>$ECN）和连续递减保护机制（$D_{max}=3$），有效防止窗口过度收缩。差异化拥塞响应（丢包0.70/ECN 0.85/超时0.50）精准匹配拥塞严重程度，克服了传统"一刀切"窗口减半的缺陷。强化学习驱动的自适应参数调优基于三级RTT梯度（1.3/2.0/3.0）、奖励指数移动平均（$\eta=0.1$）和连续增长趋势动态调整乘性增加因子$\alpha \in [1.05, 1.50]$。融合控制策略结合基于窗口最大值滤波的BDP速率控制和当前窗口保守控制，实现了快速收敛和稳健增长。

在ns-3平台9种典型数据中心网络场景的大量实验中，Lark保持与NewReno/CUBIC相当的吞吐量（差距$<$1\%，带宽利用率$>$99\%），同时将延迟降低70\%$\sim$81\%（从1.9$\sim$2.8ms降至0.4$\sim$0.8ms）。在跨Pod场景带宽利用率超过99\%，综合性能显著优于BBR。收敛速度最快（0.8秒达90\%），稳态波动最小（$\pm$3\%），公平性良好（Jain指数0.98），混合流量场景中短流完成时间降低72\%。实验结果充分验证了Lark"高吞吐、低延迟"的设计目标，证明了多信号融合、差异化响应和自适应调优等核心机制的有效性，为数据中心网络拥塞控制提供了一种高效、自适应的解决方案。

%========================================
% 参考文献
%========================================

{\small
\begin{thebibliography}{25}

\bibitem{idc2023}
IDC. Worldwide Global DataSphere Forecast, 2023-2027. IDC Report, 2023.

\bibitem{datacentertraffic}
Benson T, Akella A, Maltz D A. Network traffic characteristics of data centers in the wild. In Proc. IMC, 2010: 267-280.

\bibitem{jacobson1988}
Jacobson V. Congestion avoidance and control. ACM SIGCOMM CCR, 1988, 18(4): 314-329.

\bibitem{newreno}
Floyd S, Henderson T, Gurtov A. The NewReno modification to TCP's fast recovery algorithm. RFC 6582, 2012.

\bibitem{cubic}
Ha S, Rhee I, Xu L. CUBIC: a new TCP-friendly high-speed TCP variant. ACM SIGOPS OSR, 2008, 42(5): 64-74.

\bibitem{bufferbloat}
Gettys J, Nichols K. Bufferbloat: Dark buffers in the Internet. Commun. ACM, 2012, 55(1): 57-65.

\bibitem{vegas}
Brakmo L S, Peterson L L. TCP Vegas: End to end congestion avoidance on a global Internet. IEEE JSAC, 1995, 13(8): 1465-1480.

\bibitem{bbr}
Cardwell N, Cheng Y, et al. BBR: Congestion-based congestion control. ACM Queue, 2016, 14(5): 20-53.

\bibitem{copa}
Arun V, Balakrishnan H. Copa: Practical delay-based congestion control for the Internet. In Proc. NSDI, 2018: 329-342.

\bibitem{bbrfairness}
Hock M, Bless R, Zitterbart M. Experimental evaluation of BBR congestion control. In Proc. ICNP, 2017: 1-10.

\bibitem{compound}
Tan K, Song J, et al. A compound TCP approach for high-speed and long distance networks. In Proc. INFOCOM, 2006: 1-12.

\bibitem{illinois}
Liu S, Ba\c{s}ar T, Srikant R. TCP-Illinois: A loss-and delay-based congestion control algorithm. Perf. Eval., 2008, 65(6-7): 417-440.

\bibitem{westwood}
Mascolo S, Casetti C, et al. TCP Westwood: Bandwidth estimation for enhanced transport over wireless links. In Proc. MobiCom, 2001: 287-297.

\bibitem{rfc3168}
Ramakrishnan K, Floyd S, Black D. The addition of ECN to IP. RFC 3168, 2001.

\bibitem{dctcp}
Alizadeh M, Greenberg A, et al. Data center TCP (DCTCP). ACM SIGCOMM CCR, 2010, 40(4): 63-74.

\bibitem{hpcc}
Li Y, Miao R, et al. HPCC: High precision congestion control. In Proc. SIGCOMM, 2019: 44-58.

\bibitem{rlcc}
Jay N, Rotman N, et al. A deep reinforcement learning perspective on Internet congestion control. In Proc. ICML, 2019: 3050-3059.

\bibitem{remy}
Winstein K, Balakrishnan H. TCP ex machina: Computer-generated congestion control. ACM SIGCOMM CCR, 2013, 43(4): 123-134.

\bibitem{aurora}
Jay N, Rotman N, et al. Internet congestion control via deep reinforcement learning. In Proc. NeurIPS, 2019.

\bibitem{orca}
Abbasloo S, Yen C Y, Chao H J. Classic meets modern: A pragmatic learning-based congestion control. In Proc. SIGCOMM, 2020: 632-647.

\bibitem{pcc}
Dong M, Li Q, et al. PCC: Re-architecting congestion control for consistent high performance. In Proc. NSDI, 2015: 395-408.

\bibitem{ns3}
Riley G F, Henderson T R. The ns-3 network simulator. Springer, 2010: 15-34.

\bibitem{opengym}
Gaw\l{}owicz P, Zubow A. ns-3 meets OpenAI Gym: The playground for ML in networking research. In Proc. MSWiM, 2019: 113-120.

\bibitem{tcp_congestion}
Afanasyev A, Tilley N, et al. Host-to-host congestion control for TCP. IEEE COMST, 2010, 12(3): 304-342.

\bibitem{fairness}
Jain R K, Chiu D M W, Hawe W R. A quantitative measure of fairness and discrimination. DEC Research Report, 1984.

\bibitem{meta_learning}
Finn C, Abbeel P, Levine S. Model-agnostic meta-learning for fast adaptation of deep networks. In Proc. ICML, 2017: 1126-1135.

\end{thebibliography}
}

\end{document}
